% JAC1964-20261121, /root/arithmetic_sign.  New reduction and exact source audit.
\subsection*{AR-24. A projected modified-diagonal test for every complementary class}

\begin{proposition}[AR-24]
Let $(A,L)/K$ be an abelian threefold and let $\pi_N$ be the
self-transpose Chow projector defining the complementary summand
$N\subset h^3(A)$ used above.  Let $z\in\operatorname{CH}^2(N)_\mathbf Q$.
Choose the curve $W$ and integer $M>0$ as constructed in AR-21,
so $\operatorname{Ce}(W)=Mz$ and
$[W]=MV+b c_1(L)^2$ for an effective rational cycle $V$ and $b>0$.
Write $f:C\to W\subset A$ for the normalization.

After a finite extension $K'/K$, there is a degree-one divisor $e$ on
$C_{K'}$, with rational coefficients, for which
$\operatorname{alb}_A(f_*e)=0$.  Define
\[
 \mu:C^3\longrightarrow A,\qquad
 (x_1,x_2,x_3)\longmapsto f(x_1)+f(x_2)+f(x_3).
\]
This map is generically finite and surjective; let its degree be
$\delta>0$.  If $\Delta_e\in\operatorname{CH}^2(C^3)_\mathbf Q$ is
the Gross--Schoen modified diagonal with base divisor $e$, then
\[
 (\pi_N)_*\mu_*\Delta_e=6Mz.                         \tag{AR.24a}
\]
Moreover
\[
 E_N=\frac1\delta\,\mu^*\circ(\pi_N)_*\circ\mu_*
                                                        \tag{AR.24b}
\]
is an exact self-transpose Chow idempotent on $C^3$.  With the
functorial height pairing used throughout,
\[
 \begin{split}
 (E_N)_*\Delta_e&=\frac{6M}{\delta}\,\mu^*z,\\
 h_K(z,z)
 &=\frac{\delta}{36M^2}
       h_{K'}((E_N)_*\Delta_e,(E_N)_*\Delta_e)\\
 &=\frac{\delta}{36M^2}
       h_{K'}(\Delta_e,(E_N)_*\Delta_e).
 \end{split}                                             \tag{AR.24c}
\]
All heights are divided by their number-field degree.  Thus the field
extension introduces no additional factor in (AR.24c).
\end{proposition}

\begin{proof}
First $W$ generates $A$ as an abelian variety, in the sense that it
is not contained in a translate of a proper abelian subvariety.
Otherwise let $q:A\to B$ be a nonzero abelian quotient constant on
$W$, and take an ample line bundle $H$ on $B$.  The nonzero nef class
$D=q^*c_1(H)$ would satisfy $D\cdot W=0$.  But the effective
representative from AR-21 gives
\[
 D\cdot W=M(D\cdot V)+b(D\cdot L^2)>0.
\]
Here the first term is nonnegative, and $D\cdot L^2>0$ for a nonzero
nef divisor $D$ on an abelian threefold.  For example, over the
complex numbers a nonzero semipositive translation-invariant form
for $D$ has strictly positive integral after wedging with the square
of the positive form for $L$.  This proves the assertion.

After a finite extension choose $p\in C(K')$.  The induced homomorphism
$\operatorname{Jac}(C)\to A$ is surjective by the preceding
generation assertion.  Choose a preimage of $-f(p)$ over
$\overline K$, and enlarge $K'$ so that it and a degree-zero divisor
$D_0$ representing it are defined.  This is a finite extension:
every chosen algebraic point and divisor has a finite field of
definition.  Then $e=p+D_0$ has degree one and
$\operatorname{alb}_A(f_*e)=0$.  It need not be effective or a
single point, and we use the modified diagonal defined for an
arbitrary rational degree-one divisor.

To verify the degree of $\mu$, put $W_r=W+\cdots+W$ with $r$ summands.
If $W_r\ne A$ and $\dim(W_r+W)=\dim W_r$, then all translates
$W_r+w$, $w\in W$, coincide, since they are closed integral
subvarieties of the same dimension inside the integral variety
$W_r+W$.  Hence $W-W$ lies in the stabilizer of $W_r$.  The closed
subgroup generated by $W-W$ is $A$, because $W$ generates $A$;
this would force $W_r=A$, a contradiction.  Dimension therefore
increases until it reaches three.  Consequently $W_3=A$ and the
proper dominant map $\mu:C^3\to A$ between threefolds is
generically finite.

Set $Z=f_*e\in\operatorname{CH}^3(A)_\mathbf Q$.  The definition of
the seven modified-diagonal terms, followed by proper pushforward,
gives the exact Chow identity
\[
 \mu_*\Delta_e=[3]_*[W]
            -3[2]_*[W]*Z+3[W]*Z*Z,                    \tag{AR.24d}
\]
where $*$ is the Pontryagin product.  This remains true for a
rational divisor $e$, by bilinearity in its terms.  A degree-one
zero-cycle with Albanese sum zero has Beauville components
$Z_{(0)}=[0]$ and $Z_{(1)}=0$.  Indeed its Fourier transform in
codimensions zero and one is respectively its degree and the
Poincare divisor at its Albanese sum.  Zero-cycles have no negative
Beauville grades, as their Fourier transforms would have negative
codimension.

Pontryagin multiplication adds Beauville grades, and the curve
class $[W]\in\operatorname{CH}^2(A)$ has no negative grades by
the Fourier calculation in AR-1a.  On its grade-one component,
$[n]_*$ acts as $n^3$.  Taking grade one in (AR.24d) therefore
gives
\[
 \bigl(\mu_*\Delta_e\bigr)_{(1)}
       =(27-3\cdot8+3)[W]_{(1)}=6\operatorname{Ce}(W).
\]
Since $\pi_N$ factors through $\pi_3$ and
$\operatorname{Ce}(W)=Mz\in\operatorname{CH}^2(N)$, this proves
(AR.24a).  No identification of $f_*e$ with $[0]$ in the whole
zero-cycle Chow group was used; only its first two graded parts
were needed.

For a proper generically finite morphism between smooth projective
varieties of the same dimension, graph composition gives the exact
correspondence identity
$\mu_*\circ\mu^*=\delta\,\mathrm{id}$ on the target.  Equivalently
it is the Chow projection formula, and it is valid without a
flatness assumption on $\mu$.  Together with
$\pi_N^2=\pi_N=\pi_N^t$, this proves
$E_N^2=E_N=E_N^t$.  Equation (AR.24a) proves the first equality in
(AR.24c).  Adjointness for the height pairing gives
\[
 \begin{split}
 h_{K'}((E_N)_*\Delta_e,(E_N)_*\Delta_e)
 &=\frac{36M^2}{\delta^2}
       h_{K'}(\mu^*z,\mu^*z)\\
 &=\frac{36M^2}{\delta^2}
       h_{K'}(z,\mu_*\mu^*z)
  =\frac{36M^2}{\delta}h_K(z,z).
 \end{split}
\]
Self-transposition and idempotence give the last equality in
(AR.24c).  Normalized heights are invariant under the finite
extension.  This completes the reduction.
\end{proof}

\paragraph{What the known Gross--Schoen statements provide.}
Zhang's \emph{Gross--Schoen cycles and dualising sheaves}, Invent.
Math. \textbf{179} (2010), pp.~2--3, defines $\Delta_e$ for rational
degree-one divisors; Lemma 2.3.1, p.~23, gives homological
triviality.  Height adjointness for transposed correspondences is
Lemma 2.1.1, p.~17.  Theorem 1.3.1, p.~5, computes the height of
the whole modified diagonal from the admissible dualising
self-intersection, the Neron--Tate base-point term, and local
$\varphi$-terms.  Theorem 2.3.5, p.~27, computes pairings with
tensor products of curve correspondences as triple arithmetic
intersections on $C\times C$.

The sign conclusions needed for arbitrary projections are not
theorems of that source.  It explicitly states their semipositivity
as Conjecture 2.4.1, p.~28.  Conjecture 2.4.2 on the same page
asks for the nonpositive triple Poincare intersection for three
semipositive symmetric bundles on an abelian variety and a curve
through zero.  This is the same kind of cubic sign isolated in
AR-22; the literal through-zero hypothesis and the chosen
normalizations must still be respected when comparing the
statements.  Theorem 1.3.1 and Theorem 2.3.5 are identities and do
not settle that sign.  We have not assumed that $E_N$ is a single
tensor product of curve correspondences, nor that its image is
primitive for an arbitrary ample polarization on $C^3$.

Even a separately proved positive value for the full $\Delta_e$
would not suffice for this projection.  The exact self-adjoint
decomposition gives
\[
 h(\Delta_e,\Delta_e)
 =h(E_N\Delta_e,E_N\Delta_e)
  +h((1-E_N)\Delta_e,(1-E_N)\Delta_e),                 \tag{AR.24e}
\]
and positivity of the left side does not separately order the two
summands.  Positivity of the selected summand is precisely the
original height problem by (AR.24c).

\paragraph{The generic height theorem does not close this projection.}
The current primary Gao--Zhang manuscript, \emph{Heights of Ceresa
and Gross--Schoen cycles}, Theorem 1.1, pp.~3--4, proves a lower
bound with an additive constant on a specified nonempty open part
of moduli.  Corollary 1.2 gives its Northcott consequence.  These
apply to the whole Gross--Schoen height, not to every Chow
projector.  Our construction does not place $C$ in that open set;
even if it did, (AR.24e) would still require control of the other
summand.  In the fixed linear systems used in AR-21 the pushed
class $\operatorname{Ce}(W)=Mz$ is constant in the Chow group of
$A$, although the auxiliary curves and their other motivic
components vary.  A lower bound for their total heights cannot
isolate the sign of this constant component.

The source itself separates the metric construction from positivity:
\S9.3, p.~29, Conjecture 9.6 proposes a nef adelic extension of the
general cycle-height line.  Its final sentence states that for the
Gross--Schoen and Ceresa families the extension and volume identity
are proved, but the nefness is not.  Thus that construction cannot
be used as an already proved nonnegative-degree input for the
height line in AR-10.

The remaining gate on this route is consequently an inequality
for the explicit projected cycle $E_N\Delta_e$, or an equivalent
bound on the complementary term in (AR.24e).  The present argument
constructs that projection and checks its exact positive scaling
factor; it does not infer its positivity from a full-cycle estimate.
Primary sources:
\url{https://web.math.princeton.edu/\string~shouwu/publications/gross-schoen.pdf}
and \url{https://ziyangjeremygao.github.io/articles/HP.pdf}.
